In bibliometrics, many indicators are based on citations (or being cited) over specific periods, such as the 3‑year impact factor and the 5‑year impact factor \cite{Cla1}. From one viewpoint, these metrics are intended to ensure fairness in evaluation: older papers have had more opportunities to be cited, so the effect of a paper’s “age” is factored out. This perspective is important, for example, when selecting journals.  

However, it is also crucial to understand how citations change over the long term. First, it provides the rationale for setting the evaluation period that underlies the aforementioned notion of fairness. Second, capturing cumulative effects that cannot be seen in the short term requires truly long‑term observation. Moreover, to grasp the phenomenon of citation more deeply, it is desirable to go beyond merely counting citations (or being cited) and to approach the structural aspects, such as citation relationships.  

In this study, we adopt the above viewpoints and attempt to elucidate the long‑term dynamics of citations.
